\documentclass{article}
\usepackage[utf8]{inputenc}
\usepackage{amsmath, amssymb, amsthm}
\usepackage{geometry}
\geometry{a4paper, margin=1in}

\title{\textbf{A Deduced Theorem for Goursat-Type Integrals with Cube Roots}}
\author{}
\date{}

\theoremstyle{plain}
\newtheorem{theorem}{Theorem}
\newcommand{\Fcal}{\mathfrak{F}}

\begin{document}
\maketitle

\begin{theorem}[Analog for Cube Roots]
Let $R(t)$ be a cubic polynomial with distinct roots. An integral of the form 
$$ I = \int \frac{\Fcal(t)}{R(t)^{1/3}} dt \quad \text{or} \quad I = \int \frac{\Fcal(t)}{R(t)^{2/3}} dt $$
is elementary if it belongs to the family of Goursat-type pseudo-elliptic integrals[cite: 1]. Such an integral can be reduced by either a single substitution or by splitting the rational part $\Fcal(t)$ into at most two components, where each resulting integral is reducible by a single substitution.
\end{theorem}

\begin{proof}
The proof is a deduction based on the assertions presented in the source text[cite: 1]. The argument follows the structure established for the square-root case, adapted to the specific rules for cube roots as described.

\subsubsection*{1. Detection of Reducibility Type}
The text states that prescriptions analogous to those in Goursat's 1887 paper exist for the cube root case. These allow one to detect whether the integral can be solved with a single substitution or if the rational part, $\Fcal(t)$, must be split.

\subsubsection*{2. Case 1: No Splitting Required}
If the detection method indicates no splitting is needed, the text asserts that the integral can be rationalized by a single substitution. This substitution transforms the original integral into an elementary form. The resulting elementary integral will be of the type:
$$ \int R_1(u^3) du \quad \text{or} \quad \int u \cdot R_2(u^3) du $$
where $R_1$ and $R_2$ are rational functions. These are integrable using standard methods.

\subsubsection*{3. Case 2: Splitting Required}
If the detection method indicates that splitting is necessary, the text states that $\Fcal(t)$ must be decomposed into two components:
$$ \Fcal(t) = \Fcal_1(t) + \Fcal_2(t) $$
The text explicitly notes that "three components are never needed," which is a key distinction from the square root case involving the Klein four-group of involutions. The original integral is thus broken into two separate integrals:
$$ I = \int \frac{\Fcal_1(t)}{R(t)^{n/3}} dt + \int \frac{\Fcal_2(t)}{R(t)^{n/3}} dt \quad (n=1 \text{ or } 2)$$
Each of these two component integrals then falls under Case 1, where each succumbs to a single rationalizing substitution.

\subsubsection*{Conclusion}
Since any integral of this type can be reduced to elementary forms—either directly (Case 1) or via a two-part decomposition where each part is reducible (Case 2)—the theorem is established as described by the source text.
\end{proof}

\end{document}
